\chapter{Introduction}

Named Entity Recognition (NER) is the task of finding spans of text that
refer to real-world entities and assigning each span a type such as
\emph{Person}, \emph{Organization} or \emph{Location}. It is a basic step
in information extraction, question answering, search and knowledge-graph
construction. English NER is a mature area, but most Indian languages are
still low-resource. This project builds and analyses a NER system for
Malayalam, the language of Kerala, using a fine-tuned multilingual BERT
model.

\section{NER as Sequence Labelling}

NER is treated as token-level sequence labelling: for a sentence
$\mathbf{x} = (x_1, \ldots, x_n)$ the model predicts one tag $y_i$ per word.

\begin{definition}[IOB2 scheme]\label{def:iob2}
Each word is tagged \texttt{B-}$T$ (first word of an entity of type $T$),
\texttt{I-}$T$ (a later word inside the entity) or \texttt{O} (not an
entity)~\cite{ramshaw1995text}. With $T \in \{\mathrm{PER}, \mathrm{ORG},
\mathrm{LOC}\}$ there are seven tags. A tag \texttt{I-}$T$ is valid only
after \texttt{B-}$T$ or \texttt{I-}$T$.
\end{definition}

\section{Challenges in Malayalam}

\begin{itemize}
  \item \textbf{No capitalization.} The script has no upper case, so the
  strongest cue used in English NER is missing.
  \item \textbf{Agglutination.} Case markers fuse onto names, so one entity
  has many surface forms. For example, \emph{Thiruvananthapuram} appears as
  \emph{Thiruvananthapurathth} (``in Thiruvananthapuram'') or
  \emph{Thiruvananthapurathekk} (``to Thiruvananthapuram'').
  \item \textbf{Ambiguity.} Many names are also common words, and
  organization names often contain place names (e.g.\ \emph{Kerala
  Blasters}).
  \item \textbf{Little annotated data.} Before Naamapadam~\cite{mhaske2023naamapadam},
  systems were trained on a few hundred hand-labelled
  sentences~\cite{lakshmi2016fuzzy, sp2020towards}.
\end{itemize}

Pre-trained multilingual Transformers such as mBERT~\cite{devlin2019bert},
MuRIL~\cite{khanuja2021muril} and IndicBERT~\cite{kakwani2020indicnlpsuite,
doddapaneni2023towards} address these issues. They transfer knowledge across
languages, and their subword tokenizers split long agglutinated words into
reusable pieces.

\section{Problem Statement and Objectives}

Given a raw Malayalam sentence, identify every PER, ORG and LOC mention and
output a valid IOB2 tag sequence. The objectives are:
\begin{enumerate}
  \item Study and critically analyse the base paper~\cite{mohan2023finetuned}
  and related Malayalam NER work.
  \item Fine-tune a multilingual BERT encoder (IndicBERTv2) for Malayalam NER
  with subword-aware label alignment.
  \item Extend the base approach with a CRF output layer and compare it with
  the softmax baseline using strict entity-level metrics.
  \item Deploy the model as a web application and a REST API.
\end{enumerate}
